% TOP_PROBLEM: {"id": 30007363, "title": "One root-independent IID factor for all free uniform spanning forests", "category_id": 19, "category": {"id": 19, "name": "probability", "display_name": "Probability", "description": "Problems involving probability theory, stochastic processes, and random structures.", "slug": "probability", "order_index": 19, "created_at": "2026-07-31T15:26:25.671Z"}, "status": "open"}
% FIRST_POSTED: 2026-10-06
% ATTEMPT_STATUS: solved
% OPENAI_CLAIM: {"schema_version":1,"source_repo":"https://github.com/openai/math","source_commit":"adc7f1241b42e322a6451854ab7e4b4c146bf78a","release_date":"2026-10-06","match":"full","resolution":"proved","independently_reviewed":false,"second_pass":"agreed","families":[{"family":"231","title":"The free uniform spanning forest is a factor of IID","subject":"Probability and statistical mechanics","manuscripts":[{"title":"The free uniform spanning forest is a factor of IID","pdf_path":"preprints/The-free-uniform-spanning-forest-is-a-factor-of-IID-September-25-2026/The-free-uniform-spanning-forest-is-a-factor-of-IID-September-25-2026.pdf","readme_path":"preprints/The-free-uniform-spanning-forest-is-a-factor-of-IID-September-25-2026/README.md","date":"2026-09-25","theorem_ref":"thm:main (sections/introduction.tex:54); thm:sr-group (sections/introduction.tex:102)","statement_sources":[{"path":"preprints/The-free-uniform-spanning-forest-is-a-factor-of-IID-September-25-2026/build/sections/introduction.tex","line":54,"label":"thm:main"},{"path":"preprints/The-free-uniform-spanning-forest-is-a-factor-of-IID-September-25-2026/build/sections/introduction.tex","line":102,"label":"thm:sr-group"}]}],"lean":{"doc_path":"lean/docs/231.md","comparators":["lean/ComparatorChallenges/FreeUniformSpanningForest.json","lean/ComparatorChallenges/FreeUniformSpanningForest.lean","lean/ComparatorChallenges/StronglyRayleighDPP.json","lean/ComparatorChallenges/StronglyRayleighDPP.lean"],"declarations":[],"covers_main_theorem":true},"reasoning_trace":null}]}
\documentclass[11pt]{article}
\usepackage[margin=25mm]{geometry}
\usepackage{fontspec}
\setmainfont{Latin Modern Roman}
\usepackage{amsmath,amssymb,mathtools}
\usepackage{unicode-math}
\setmathfont{Latin Modern Math}
\usepackage{enumitem}
\usepackage{xurl,hyperref,bookmark}
\hypersetup{colorlinks=true,urlcolor=blue}
\setcounter{secnumdepth}{0}
\setlength{\parindent}{0pt}
\setlength{\parskip}{5pt}
\setlength{\emergencystretch}{3em}
\begin{document}
\section*{One root-independent IID factor for all free uniform spanning forests}

\subsection{OpenAI's claim}
OpenAI presents a proof in the selected manuscript. The free spanning forest on every infinite locally finite simple graph is claimed obtainable by one root-free Borel rule from iid labels; invariant strongly Rayleigh processes have an analogous group-equivariant construction.

\subsection{Mathematical statement of the claimed result}
\paragraph{Result 1.} One Borel root-independent rule \(\Phi(G,U,e)\in\{0,1\}\) is claimed to work for every infinite graph that is connected, has finite degrees, and is simple, undirected and unweighted. With independent uniform vertex labels, its selected edge set
\[
 \{e:\Phi(G,U,e)=1\}
\]
has the free uniform spanning-forest law and the rule commutes with graph isomorphisms. The second theorem takes any countable group \(\Gamma\) and any translation-invariant strongly Rayleigh law \(\nu\) on \(\{0,1\}^\Gamma\). It supplies a Borel map \(\Phi_\nu:[0,1]^\Gamma\to\{0,1\}^\Gamma\), equivariant at every input, whose iid-uniform input has output law \(\nu\).

\subsection{Outline of the argument}
Finite determinant tilts define a canonical response drift for the free forest law. Brownian observations provide independent innovations, and a deterministic inverse reconstructs the sampler; vertex labels are converted to independent edge paths in a measurable rule that also applies to varying graphs.

\subsection{Scope relative to this problem}
The free spanning forest on every infinite locally finite simple graph is claimed obtainable by one root-free Borel rule from iid labels; invariant strongly Rayleigh processes have an analogous group-equivariant construction. The question retains these same parameters and conventions. Within these assumptions, this is a claimed affirmative answer to the stated question.

\subsection{Formal verification}
For family 231, the Lean scope includes a root-free iid rule for free spanning forests and invariant strongly Rayleigh and determinantal sampling supplements without amenability assumptions. The listed Comparator files are \texttt{FreeUniformSpanningForest.lean}, \texttt{StronglyRayleighDPP.lean}. The inventory does not supply a named declaration for these Comparator entries. This site has not executed the Lean code. The scope descriptions do not certify the remaining manuscript arguments.

\subsection{Status caveat}
OpenAI attributes these manuscripts to work by an internal model. Its release cautions that arguments without formalization can contain mistakes. This catalogue has not independently checked the proofs; the labels report the mathematical claims attributed to OpenAI.

\subsection{Sources}
\begin{itemize}
\item OpenAI, \emph{The free uniform spanning forest is a factor of IID}, 2026-09-25; citation key \texttt{OAI:The-free-uniform-spanning-forest-is-a-factor-of-IID-September-25-2026}. \href{https://github.com/openai/math/blob/adc7f1241b42e322a6451854ab7e4b4c146bf78a/preprints/The-free-uniform-spanning-forest-is-a-factor-of-IID-September-25-2026/The-free-uniform-spanning-forest-is-a-factor-of-IID-September-25-2026.pdf}{PDF (GitHub)}; \href{https://github.com/openai/math/raw/adc7f1241b42e322a6451854ab7e4b4c146bf78a/preprints/The-free-uniform-spanning-forest-is-a-factor-of-IID-September-25-2026/The-free-uniform-spanning-forest-is-a-factor-of-IID-September-25-2026.pdf}{PDF (direct)}; \href{https://github.com/openai/math/blob/adc7f1241b42e322a6451854ab7e4b4c146bf78a/preprints/The-free-uniform-spanning-forest-is-a-factor-of-IID-September-25-2026/README.md}{Manuscript page and citation}; \href{https://github.com/openai/math/blob/main/preprints/The-free-uniform-spanning-forest-is-a-factor-of-IID-September-25-2026/The-free-uniform-spanning-forest-is-a-factor-of-IID-September-25-2026.pdf}{Latest version}; Statement source: \href{https://github.com/openai/math/blob/adc7f1241b42e322a6451854ab7e4b4c146bf78a/preprints/The-free-uniform-spanning-forest-is-a-factor-of-IID-September-25-2026/build/sections/introduction.tex\#L54}{thm:main}; Statement source: \href{https://github.com/openai/math/blob/adc7f1241b42e322a6451854ab7e4b4c146bf78a/preprints/The-free-uniform-spanning-forest-is-a-factor-of-IID-September-25-2026/build/sections/introduction.tex\#L102}{thm:sr-group}.
\item Family 231: \href{https://github.com/openai/math/blob/adc7f1241b42e322a6451854ab7e4b4c146bf78a/lean/docs/231.md}{Lean scope}.
\item \href{https://github.com/openai/math/blob/adc7f1241b42e322a6451854ab7e4b4c146bf78a/lean/ComparatorChallenges/FreeUniformSpanningForest.json}{Comparator FreeUniformSpanningForest.json}.
\item \href{https://github.com/openai/math/blob/adc7f1241b42e322a6451854ab7e4b4c146bf78a/lean/ComparatorChallenges/FreeUniformSpanningForest.lean}{Comparator FreeUniformSpanningForest.lean}.
\item \href{https://github.com/openai/math/blob/adc7f1241b42e322a6451854ab7e4b4c146bf78a/lean/ComparatorChallenges/StronglyRayleighDPP.json}{Comparator StronglyRayleighDPP.json}.
\item \href{https://github.com/openai/math/blob/adc7f1241b42e322a6451854ab7e4b4c146bf78a/lean/ComparatorChallenges/StronglyRayleighDPP.lean}{Comparator StronglyRayleighDPP.lean}.
\item \href{https://github.com/openai/math/blob/adc7f1241b42e322a6451854ab7e4b4c146bf78a/CONTENTS.md}{OpenAI catalogue entry}.
\end{itemize}
\end{document}
